% Generated by task tex-boundary from dossiers/boundary-adaptor.md, sections 0 and A (A.1 to A.6).
% Breakable inline code for this fragment's tables and prose (LuaLaTeX): \bndc{...} sets its
% argument verbatim in the monospace font, with break points after punctuation and between the
% words of a camel-case identifier. Defined once, whichever boundary fragment is read first.
\ifdefined\bndc\else
\directlua{
  bndbreak = function(s)
    local bs = string.char(92)
    local hash = string.char(35)
    s = s:gsub(hash .. hash, hash)
    local after = { ["."] = true, ["_"] = true, ["/"] = true, [":"] = true, [","] = true,
      ["-"] = true, [")"] = true, ["]"] = true, ["="] = true, ["|"] = true }
    local prev = ""
    for _, c in utf8.codes(s) do
      local ch = utf8.char(c)
      if ch:match("^[A-Z]$") and prev:match("^[a-z]$") then
        tex.sprint(bs .. "penalty400{}")
      end
      tex.sprint(-2, ch)
      if after[ch] or c > 127 then
        tex.sprint(bs .. "allowbreak{}")
      end
      prev = ch
    end
  end
}
\newcommand{\bndc}[1]{\texttt{\directlua{bndbreak("\luaescapestring{\detokenize{#1}}")}}}
\fi

\section{The adaptor: the two relations side by side}\label{sec:gen-boundary-relations}

\begin{sloppypar}
Short forms in this section. leanerVM is at \code{b435631}; \code{bp:} is the blueprint
\file{docs/roadmap/protocol-blueprint.md} there. Lean files are named by their last component:
\file{Statement.lean}, \file{Boundary.lean}, \file{Channels.lean}, \file{Bytecode.lean} are under
\file{LeanerVM/Arithmetization/}; \file{Memory.lean}, \file{Execution.lean} under
\file{LeanerVM/Semantics/}; \file{Spine/Instance.lean}, \file{Spine/Toy.lean},
\file{FixedColumns.lean}, \file{Stack.lean}, \file{ToArkLib/Refinement.lean} under
\file{LeanerVM/Protocol/}; all at \code{b435631}. The arithmetization, semantics and parameters
modules cited are byte-identical between \code{b435631} and \code{144c5aa}, and
\file{Spine/Instance.lean} differs by one import line. Clean is cited at its old pin
\code{93c9d1ef} (\file{FlatEnsemble.lean}, \file{FlatComponent.lean} under \file{Clean/Air/},
\file{Operations.lean}, \file{Expression.lean} under \file{Clean/Circuit/}), ArkLib at
\code{dca90385}, CompPoly at \code{3468b38c}. Rust files are under \file{crates/lean_vm/src/}
(\code{layout.rs} is \file{cpu/layout.rs}), at leanVM's pin \code{a386121f}, read with
\code{git show a386121f:<path>}. The number after a colon is the line.
\end{sloppypar}

The proof system proves knowledge of a stack satisfying \bndc{M3Holds I input q}, a relation on
one committed column (\bndc{Spine/Instance.lean:219-220}); the arithmetization consumes
\bndc{SatisfiedBy prog input w}, a relation on a Clean \bndc{EnsembleWitness}
(\bndc{Statement.lean:311-340}). The \emph{adaptor} between them is sketched in the blueprint's
Layers 2 and 3 (\bndc{bp:757-867}) and is not built. This section sets the two relations side by
side, clause by clause, and asks what the adaptor's two theorems need.

\subsection{The objects}\label{sec:gen-boundary-objects}

Every object below is quoted from the pinned sources.
\begin{itemize}
\item \textbf{Clean} (\bndc{93c9d1ef}) is the circuit library the arithmetization is written in.
  A \bndc{Component F} is one AIR table's row circuit: an input row type, and a
  \bndc{GeneralFormalCircuit} with its constraints, its bus interactions and its specification.
  An \bndc{Ensemble F PublicIO} is a list of components, a list of bus channels and one
  distinguished \enquote{verifier} component read once on the public input
  (\bndc{FlatEnsemble.lean:11-17}). An \bndc{EnsembleWitness ens} is the prover's data for it:
  one \bndc{Table} (a list of raw rows, \bndc{Array F}) per component, the string-keyed
  \bndc{ProverData} (a function \bndc{String → (n : ℕ) → Array (Vector F n)},
  \bndc{Expression.lean:21-22}), and the public input (\bndc{FlatEnsemble.lean:19-25}). It lives
  in \bndc{Type 1} because a \bndc{Component} bundles type-level fields (probe \bndc{Shapes},
  output lines 1--4). \bndc{w.Constraints} is \enquote{every component's asserted expressions
  vanish on every row and every lookup is contained} (\bndc{Operations.lean:687-688}; bus
  guarantees are \emph{not} part of it, probe \bndc{Shapes}, output lines 154--159), and
  \bndc{w.interactions} is the concatenation of every table's emitted bus interactions, the
  verifier's one row first (\bndc{FlatEnsemble.lean:225-226}).
\item \textbf{The spine} (\bndc{LeanerVM/Protocol/Spine/Instance.lean}) is the proof system's
  abstract view of an arithmetization. An \bndc{M3Instance} (\bndc{:119-148}) is: a \bndc{Shape}
  (a number of tables, each with a log-height \bndc{τ} and a width); a statement type
  \bndc{Stmt}; per table, a list of constraint polynomials and a list of flush tuples (a
  \bndc{Side}, push or pull, and sixteen coordinate polynomials, the domain separator first); a
  degree bound \bndc{d} with its two proofs; the count columns; the boundary blocks; the
  log-height \bndc{μ} of the one committed column and a \bndc{Layout} reading every column off
  it; the public lines; and the Flock predicate \bndc{aux}. A \bndc{Coord} of a boundary block
  (\bndc{:86-89}) is a constant, a \bndc{known} column (data both parties hold: the index column,
  the program) or a \bndc{committed} column of a table. \bndc{M3Holds I input q}
  (\bndc{:219-220}) is the conjunction of five clauses on the column \bndc{q}.
\item \textbf{ArkLib} (\bndc{dca90385}) is the proof-system library. Its knowledge-soundness game
  \bndc{Verifier.knowledgeSoundnessWith} (\bndc{Security/Basic.lean:299-323}) fixes a statement,
  runs a malicious prover against the verifier, applies a named straight-line extractor to the
  transcript, and bounds the probability that the verifier accepts \emph{and} no witness in the
  extractor's \bndc{Option} slot is valid. Statement and witness types are in \bndc{Type}. A
  \bndc{Refinement R S} (\bndc{ToArkLib/Refinement.lean:31-36}) is a map on witnesses sending
  every witness valid for \bndc{R} to one valid for \bndc{S}; its two witness universes are
  independent (probe \bndc{Shapes}, output lines 10--13).
\item \textbf{CompPoly} (\bndc{3468b38c}) supplies the computable polynomials:
  \bndc{CMvPolynomial n K} (sparse maps from monomials to coefficients, with \bndc{eval} and
  \bndc{totalDegree}) and the hypercube tables \bndc{CMlPolynomialEval}. Mathlib's
  \bndc{MvPolynomial} is the classical object proofs are easiest in; \bndc{fromCMvPolynomial}
  maps the computable one to it (\bndc{MvPolyEquiv/Core.lean:35}), and \bndc{toCMvPolynomial}
  maps back and is \bndc{noncomputable} (\bndc{:41}).
\end{itemize}

\subsection{The relation leanISA consumes: \texorpdfstring{\protect\texttt{SatisfiedBy}}{SatisfiedBy}}\label{sec:gen-boundary-satisfiedby}

\begin{leancode}
structure SatisfiedBy (prog : Program) (input : PublicInput)
    (w : EnsembleWitness (leanIsaEnsemble prog)) : Prop where
  public_input_eq : w.publicInput = PublicIO.ofInput input
  constraints : w.Constraints
  state_balanced : BalancedPair w StatePull.toRaw StatePush.toRaw
  mem_balanced : BalancedPair w MemPull.toRaw MemPush.toRaw
  bytecode_balanced : BalancedPair w BytecodePull.toRaw BytecodePush.toRaw
  counts_nonzero : CountsNonzero w
  caps : Caps w
  index_columns : IndexColumnsAreRowIndices w
  seed_rows : SeedRowsAreTheImage w
  bytecode_rows : BytecodeRowsAreTheProgram prog w
  blake2s_valid : Blake2sRowsValid w
  word0_eq : (imageOf w.data).2.read (gpow 0) = some input.word0
  word1_eq : (imageOf w.data).2.read (gpow 1) = some input.word1
\end{leancode}
\src{LeanerVM/Arithmetization/Statement.lean:311-340 at b435631, docstrings omitted (probe Shapes prints it, output lines 29--55)}

The thirteen conjuncts, as the definitions below unfold them:
\begin{sloppypar}
\begin{itemize}
\item \bndc{public_input_eq}: the witness's public input is the four lanes of \bndc{input}; no
  component reads them (\bndc{Boundary.lean:120-123}).
\item \bndc{constraints}: Clean's \bndc{w.Constraints}, for every table including the verifier's
  one-row table \bndc{(∀ e ∈ ops.constraints, env e = 0) ∧ (∀ l ∈ ops.lookups, l.Contains env)}
  (\bndc{FlatEnsemble.lean:216}, \bndc{FlatComponent.lean:184-186},
  \bndc{Operations.lean:687-688}). leanISA emits no lookup (a grep of
  \bndc{LeanerVM/Arithmetization/Tables/} finds the word in prose only) and only \bndc{jumpTable}
  asserts (\bndc{Tables/Jump.lean:197-198}, the two residuals on its two local witness cells
  \bndc{w}, \bndc{b}, \bndc{:195-196}).
\item \bndc{state_balanced}, \bndc{mem_balanced}, \bndc{bytecode_balanced}: for each channel
  pair, the messages sent on the channel of that name, multiplicities ignored, form a
  permutation of each other, counted in \bndc{ℕ} (\bndc{BalancedPair}).
\item \bndc{counts_nonzero}: message index 1 of every pull on the memory and bytecode channels,
  in the six opcode tables, is nonzero (the count product).
\item \bndc{caps}: three of \bndc{read_public}'s eight checks, plus a shape of the data.
\item \bndc{index_columns}, \bndc{seed_rows}, \bndc{bytecode_rows}: the three facts Clean cannot
  express. The memory block's index column is \bndc{gpow i}; the memory block's rows are the
  image the prover data names; the bytecode block's rows are \bndc{bytecodeRowOf prog i}.
\item \bndc{blake2s_valid}: Flock's conjunct; every \bndc{BLAKE2S} row satisfies the compression
  relation.
\item \bndc{word0_eq}, \bndc{word1_eq}: the \emph{whole} \bndc{E}-word of the image at index
  \bndc{gpow 0} (resp.\ \bndc{gpow 1}) is \bndc{input.word0} (resp.\ \bndc{input.word1}), top limb
  included.
\end{itemize}
\end{sloppypar}

One level down, in the same file:
\begin{leancode}
-- :174-181  the ensemble: eight components, six channels, the program-bearing verifier
def leanIsaEnsemble (prog : Program) : Ensemble K PublicIO where
  tables := [⟨xorTable⟩, ⟨mulTable⟩, ⟨setTable⟩, ⟨derefTable⟩, ⟨jumpTable⟩, ⟨blake2sTable⟩,
    ⟨memTable⟩, ⟨bytecodeTable⟩]
  channels := [StatePull.toRaw, StatePush.toRaw, MemPull.toRaw, MemPush.toRaw,
    BytecodePull.toRaw, BytecodePush.toRaw]
  verifier := leanIsaVerifier prog
-- :207-209  messages by channel NAME, multiplicities ignored
def messagesOn (w : EnsembleWitness (leanIsaEnsemble prog)) (c : RawChannel K) :
    List (Array K) :=
  (w.interactions.filter (·.channel.name = c.name)).map (·.msg)
-- :229-231  balance: a permutation of message lists, counted in ℕ
def BalancedPair (w : EnsembleWitness (leanIsaEnsemble prog)) (pull push : RawChannel K) :
    Prop :=
  (messagesOn w push).Perm (messagesOn w pull)
-- :240-243  the count product: message index 1 of every pull of the six tables
def CountsNonzero (w : EnsembleWitness (leanIsaEnsemble prog)) : Prop :=
  ∀ t ∈ w.tables.take 6, ∀ i ∈ t.interactions,
    (i.channel.name = MemPull.name ∨ i.channel.name = BytecodePull.name) →
      ∀ h : 1 < i.msg.size, i.msg[1] ≠ 0
-- :250-263  three of read_public's eight checks, plus a shape of the data
structure Caps (w : EnsembleWitness (leanIsaEnsemble prog)) : Prop where
  minLogMem_le : minLogMem ≤ (imageOf w.data).1
  le_maxLogMem : (imageOf w.data).1 ≤ maxLogMem
  heights : ∀ t ∈ w.tables, ∃ τ ≤ maxLogRows, t.table.length = 2 ^ τ
  blake2s_height : 2 ^ minLogRowsBlake2s ≤ (blake2sRows w).length
  well_shaped : WellShapedData w.data
-- :272-273, :279-283, :290-293  the three facts Clean cannot express
def IndexColumnsAreRowIndices (w : EnsembleWitness (leanIsaEnsemble prog)) : Prop :=
  ∀ i (hi : i < (memBlockRows w).length), (memRowAt w (memBlockRows w)[i]).idx = gpow i
def SeedRowsAreTheImage (w : EnsembleWitness (leanIsaEnsemble prog)) : Prop :=
  ∃ idx cntFin : Fin (2 ^ (imageOf w.data).1) → K,
    memBlockRows w = List.ofFn fun i ↦
      (toElements (⟨idx i, cntFin i, #v[((imageOf w.data).2 i).limb 0,
        ((imageOf w.data).2 i).limb 1, ((imageOf w.data).2 i).limb 2]⟩ : MemRow K)).toArray
def BytecodeRowsAreTheProgram (prog : Program) (w : EnsembleWitness (leanIsaEnsemble prog)) :
    Prop :=
  ∃ cntFin : Fin (2 ^ prog.logSize) → K,
    bytecodeBlockRows w = List.ofFn fun i ↦ (toElements (bytecodeRowOf prog i (cntFin i))).toArray
-- :301-302  Flock's conjunct
def Blake2sRowsValid (w : EnsembleWitness (leanIsaEnsemble prog)) : Prop :=
  ∀ row ∈ blake2sRows w, Blake2sRelation (blake2sRowAt w row)
\end{leancode}
\src{LeanerVM/Arithmetization/Statement.lean at b435631, lines as marked; docstrings and the ensemble's proof field omitted; the comments are the review's}

and, in \file{Channels.lean}, the image read off the prover data:
\begin{leancode}
def imageOf (data : ProverData K) : (κ : ℕ) × MemImage κ :=
  ⟨min (Nat.log 2 (memRows data).size) maxLogMem,
    fun i ↦ (((memRows data)[(i : ℕ)]?).map fun v ↦ E.ofLimbs v[0] v[1] v[2]).getD 0⟩
\end{leancode}
\src{LeanerVM/Arithmetization/Channels.lean:188-190 at b435631}

In \file{Boundary.lean} are the three block components and their honest rows: \bndc{memTable}
(\bndc{:157-160}, push \bndc{(idx, 1, m)}, pull \bndc{(idx, cntFin, m)}), \bndc{bytecodeTable}
(\bndc{:240-243}, push \bndc{(idx, 1, opcode, op)}, pull \bndc{(idx, cntFin, opcode, op)}),
\bndc{leanIsaVerifier prog} (\bndc{:311-314}, push \bndc{(1, 1)}, pull
\bndc{(const prog.finalPc, 1)}), \bndc{memRowOf} (\bndc{:184-185}) and \bndc{bytecodeRowOf}
(\bndc{:274-276}: \bndc{⟨gpow i, cntFin, e[0], #v[e[1], …, e[7]]⟩} with
\bndc{e := entry (prog.code i)}).

\subsection{The relation the proof system proves: \texorpdfstring{\protect\texttt{M3Holds}}{M3Holds}}\label{sec:gen-boundary-m3holds}

\begin{leancode}
-- :197-198
def ConstraintsVanish (q : Column I.μ) : Prop :=
  ∀ j, ∀ C ∈ I.constraints j, ∀ x, C.eval (I.row q j x) = 0
-- :201
def Balanced (q : Column I.μ) : Prop := (I.tuples q .push).Perm (I.tuples q .pull)
-- :204-205
def CountsNonzero (q : Column I.μ) : Prop :=
  ∀ j, ∀ i ∈ I.counts j, ∀ x, (I.column q ⟨j, i⟩).values.get x ≠ 0
-- :208-211
def PublicLinesHold (input : I.Stmt) (q : Column I.μ) : Prop :=
  ∀ l ∈ I.publicLines input,
    (I.column q l.col).values.get ⟨0, Nat.two_pow_pos _⟩ = l.cell0 ∧
      (I.column q l.col).values.get ⟨1, Nat.one_lt_two_pow l.pos.ne'⟩ = l.cell1
-- :219-220
def M3Holds (I : M3Instance) (input : I.Stmt) (q : Column I.μ) : Prop :=
  I.ConstraintsVanish q ∧ I.Balanced q ∧ I.CountsNonzero q ∧ I.PublicLinesHold input q ∧ I.aux q
\end{leancode}
\src{LeanerVM/Protocol/Spine/Instance.lean at b435631, lines as marked, docstrings omitted (probe Shapes prints M3Holds, output lines 100--134)}

The five clauses are: every constraint polynomial vanishes on every row of its table; the pushed
tuples are a permutation of the pulled tuples; every cell of every count column is nonzero; cells
0 and 1 of every column a public line names hold the line's two values; and the instance's
predicate \bndc{aux} holds. The tuples of a side are the tables' flushes evaluated on every row,
then the boundary blocks' tuples, coordinate by coordinate:
\begin{leancode}
def coordCell (q : Column I.μ) {κ : ℕ} : Coord I.toShape κ → Fin (2 ^ κ) → K
  | .const c, _ => c
  | .known col, x => col.values.get x
  | .committed c h, x => (I.column q c).values.get (Fin.cast (congrArg (2 ^ ·) h.symm) x)
def tuples (q : Column I.μ) (s : Side) : List (Vector K 16) :=
  I.flushTuples q s ++ I.boundaryTuples q s
\end{leancode}
\src{LeanerVM/Protocol/Spine/Instance.lean:167-186 at b435631, docstrings and the definitions of flushTuples and boundaryTuples omitted}

A \bndc{known} coordinate is read from data the instance holds, a \bndc{const} coordinate is a
constant, and only a \bndc{committed} coordinate is read from \bndc{q}.

\subsection{The leanISA instance the blueprint implies}\label{sec:gen-boundary-instance}

Read from Layer 3 (\bndc{bp:794-867}), the status file's decision that \enquote{the fixed
columns are \bndc{Coord.known} data} (decision 8, \bndc{protocol-status.md:220-223}) and the
ground truth (\bndc{layout.rs:352-395}; the bus dossier's table of the boundary blocks):
\begin{itemize}
\item \textbf{Tables 0--5}: the six opcode components through \bndc{Component.toM3} (Layer 2):
  width \bndc{component.width} (\bndc{JUMP}'s includes its two local witness cells),
  \bndc{τ j := s.τ j}, the constraint polynomials, the flush tuples with separator
  \bndc{g^0/g^1/g^2} first and side from \bndc{channelDir}, the count columns.
\item \textbf{Column groups}: the six shared committed columns, \enquote{tables of the instance
  with no constraints and no flushes} (\bndc{bp:836-839}). They cannot be one table: their heights
  differ (\bndc{mem_0, mem_1, mem_2, cntfin_mem} at \bndc{2^logMem}; \bndc{cntfin_bc} at
  \bndc{2^prog.logSize}, \bndc{layout.rs:150}; \bndc{q_flock} at \bndc{2^(τ_5 + 8)},
  \bndc{layout.rs:155} with \bndc{K_LOG = 14}, \bndc{LOG_PACKING = 6}). So at least three tables
  with empty \bndc{constraints}, \bndc{flushes} and \bndc{counts}, which the spine's \bndc{Shape}
  (\bndc{Spine/Instance.lean:57-63}) and \bndc{M3Instance.counts} allow.
\item \textbf{Six boundary blocks}: \bndc{Coord.known (idxColumn κ)} for the index,
  \bndc{Coord.known (bytecodeSlotColumn prog s)} for the eight program columns
  (\bndc{FixedColumns.lean:75-76}), \bndc{Coord.committed} for the limbs and the finalize counts,
  \bndc{Coord.const} for the separators, the seed count \bndc{1} and the state boundary
  (\bndc{κ = 0}). The state pull's counter is \bndc{const prog.finalPc}, as
  \bndc{Boundary.lean:314} and \bndc{layout.rs:357}.
\item \textbf{Three public lines} on \bndc{mem_0, mem_1, mem_2} (\bndc{bp:834-836}): cells
  \bndc{(lanes 0, lanes 2)}, \bndc{(lanes 1, lanes 3)}, \bndc{(0, 0)}, the first two \bndc{sent}.
\item \textbf{\bndc{aux}}: Flock's R1CS on \bndc{q_flock} (the strong reading, decision 12),
  \bndc{decAux} its check.
\end{itemize}

\subsection{Conjunct by conjunct}\label{sec:gen-boundary-conjuncts}

\Cref{tab:gen-boundary-conjuncts} takes \bndc{w := witnessOf prog s q} and
\bndc{I := leanIsaInstance prog s}, with the hypotheses \bndc{hs : s.Admissible prog} and
\bndc{h : M3Holds I input q}, and says for each conjunct of \bndc{SatisfiedBy} what delivers it in
\bndc{satisfiedBy_witnessOf}, through which lemma, and whether that works. The earlier review's
table (\bndc{docs/reviews/protocol-spine.md:331-342}) was redone from the definitions.

\begin{landscape}
{\footnotesize\setlength{\tabcolsep}{4pt}
\begin{longtable}{@{}>{\raggedright\arraybackslash}p{0.13\linewidth}>{\raggedright\arraybackslash}p{0.25\linewidth}>{\raggedright\arraybackslash}p{0.34\linewidth}>{\raggedright\arraybackslash}p{0.24\linewidth}@{}}
\caption{The conjuncts of \code{SatisfiedBy} against the clauses of \code{M3Holds}, in the soundness direction \code{satisfiedBy_witnessOf}.}\label{tab:gen-boundary-conjuncts}\\
\toprule
Conjunct & Delivered by & Through & Verdict \\
\midrule
\endfirsthead
\toprule
Conjunct & Delivered by & Through & Verdict \\
\midrule
\endhead
\bottomrule
\endfoot
\bndc{public_input_eq} & clause 4 (\bndc{PublicLinesHold}), lines 1--2, if \bndc{witnessOf} reads the lanes off cells 0, 1 of \bndc{mem_0}, \bndc{mem_1}; or by construction if the map takes \bndc{input} (the \bndc{Refinement.map} has it, the blueprint's \bndc{witnessOf prog s q} does not) & \bndc{limb_ofLimbs} & works either way; harmless as a built part, since no component reads the lanes (\bndc{Boundary.lean:120-123}) \\
\bndc{constraints} & clause 1 (\bndc{ConstraintsVanish}) for tables 0--5; vacuous for tables 6, 7 and the verifier (no assert) & Layer 2's \bndc{toM3_constraints_iff}, whose right side must be Clean's \bndc{Operations.ConstraintsHold} (the two-conjunct form, \bndc{Operations.lean:687}), not the flat \bndc{constraintsHold_iff_forall_mem} of \bndc{Operations.lean:168-182} the blueprint cites (\bndc{bp:292}), which includes the bus guarantees & works, provided \bndc{witnessOf} writes the \bndc{JUMP} witness cells read from \bndc{q} into the rows; citation to fix \\
\bndc{state_balanced}, \bndc{mem_balanced}, \bndc{bytecode_balanced} & clause 2 (\bndc{Balanced}), one permutation of sixteen-tuples, filtered by coordinate 0 (the separator) and mapped back to messages & \bndc{List.Perm.filter}, injectivity of \bndc{busTuple} on arrays of one size; Layer 2's \bndc{toM3_flushes_eq} for tables 0--5; \textbf{a lemma the blueprint does not have} for the memory block, the bytecode block and the verifier: their messages in \bndc{w} equal the instance's boundary tuples (\cref{sec:gen-boundary-trap}) & works only with the missing lemma and only because the boundary's \bndc{known} data are \bndc{prog}'s (\cref{sec:gen-boundary-trap}) \\
\bndc{counts_nonzero} & clause 3 (\bndc{CountsNonzero}) & per table, the pull's count coordinate is \bndc{var c} for \bndc{c ∈ I.counts j} (\bndc{memRead}, \bndc{bytecodeRead}, \bndc{Channels.lean:244-254}: the count is a column in every table) & works if \bndc{Component.toM3} computes \bndc{count} as \enquote{coordinate 2 of every pull on a lookup channel, when it is a variable}; the blueprint gives \bndc{toM3} no way to know which channels are lookups (\bndc{bp:773-774}; \cref{f:boundary-witnessof-input-counts}) \\
\bndc{caps} & nothing in \bndc{M3Holds}; \bndc{hs} and the construction: tables of exactly \bndc{2^τ} rows, \bndc{"mem"} data of \bndc{2^logMem} rows & \bndc{Admissible} must contain \bndc{16 ≤ logMem ≤ 32}, \bndc{τ_j ≤ 32}, \bndc{3 ≤ τ_5}; table 7's height is \bndc{2^prog.logSize ≤ 2^32} by \bndc{Program.logSize_le} and \bndc{maxLogBytecode = maxLogRows} & sound: the verifier checks the sizes (\bndc{cpu/mod.rs:157-176}); \bndc{WellShapedData} is a shape of the built data, no check \\
\bndc{index_columns} & construction: \bndc{idx := gpow i} & matched by \bndc{Coord.known (idxColumn logMem)} in both memory blocks (\bndc{idxColumn_get}, \bndc{FixedColumns.lean:50-53}) & sound (\cref{sec:gen-boundary-trap}) \\
\bndc{seed_rows} & construction: \bndc{"mem"} table and block rows both from \bndc{q}'s \bndc{mem_0, mem_1, mem_2}, \bndc{cntFin} from \bndc{cntfin_mem} & \bndc{imageOf} returns \bndc{logMem} only if \bndc{logMem ≤ 32} (the \bndc{min} in \bndc{Channels.lean:189}): needs \bndc{hs} & sound \\
\bndc{bytecode_rows} & construction: \bndc{bytecodeRowOf prog i (cntfin_bc[i])} & matched by \bndc{Coord.known (bytecodeSlotColumn prog k)}, \bndc{k = 3..10}, \bndc{encodeSlots = entry} at slots 3--10 (\bndc{Bytecode.lean:94-96}) & sound (\cref{sec:gen-boundary-trap}); the blueprint's \bndc{bytecodeColumn} of \file{FixedColumns.lean} is the sixteen-slot table, whose eight slices \bndc{bytecodeSlotColumn} are the block's columns (\bndc{slice_bytecodeColumn}, \bndc{FixedColumns.lean:90-92}) \\
\bndc{blake2s_valid} & clause 5 (\bndc{aux}) & the Flock roadmap's lemma \enquote{the R1CS holds of the bits packed into \bndc{q_flock} $⇒$ the eighteen limb slots compress} (decision 12); the limbs of the row are read strided off \bndc{q_flock} & works only with the lemma, which has no carrier in any sketched signature (\cref{f:boundary-flock-circular}) \\
\bndc{word0_eq}, \bndc{word1_eq} & clause 4, all three lines & \bndc{MemImage.read_gpow} (\bndc{κ < 64} from \bndc{hs}), \bndc{ofLimbs_limb}; the third line supplies limb 2 = 0 & works; the third line is load-bearing (\cref{sec:gen-boundary-toplimb}) \\
\end{longtable}}
\end{landscape}

Two conjuncts hold trivially of the built witness (\bndc{index_columns}, \bndc{bytecode_rows})
and one half-built (\bndc{seed_rows}); this is where the trap of the built parts lives. Every
conjunct has a source; none is delivered by nothing once the caps hypothesis, the boundary lemma,
the count derivation and the Flock consequence lemma are in place.

\subsection{The trap of the built parts}\label{sec:gen-boundary-trap}

\paragraph{What \texorpdfstring{\protect\code{witnessOf}}{witnessOf} builds and what it reads.}
From \bndc{q} through \bndc{I.layout}: the rows of tables 0--5 (the \bndc{BLAKE2S} limbs strided
off \bndc{q_flock}), the memory limbs and the two finalize counts. From \bndc{prog}: the eight
entry cells of every bytecode-block row (\bndc{bytecodeRowOf prog i}) and the sentinel constant
of the verifier component. From arithmetic: the index column \bndc{gpow i} of the memory block.
From nothing: the \bndc{"mem"} table of the data is a copy of the limbs. So \bndc{w} is a
\textbf{mixed witness}, and \bndc{SatisfiedBy}'s \bndc{bytecode_balanced} and
\bndc{mem_balanced} are statements about lists that contain the built rows:
\bndc{messagesOn w BytecodePush} includes, for every slot \bndc{i}, the built message
\bndc{(g^i, 1, entry (prog.code i))} (\bndc{bytecodeTable.main}, \bndc{Boundary.lean:240-243}, on
the row \bndc{bytecodeRowOf prog i _}), and \bndc{messagesOn w BytecodePull} the read messages of
tables 0--5, taken from \bndc{q}.

\paragraph{What \texorpdfstring{\protect\code{M3Holds}}{M3Holds} establishes.}
\bndc{I.Balanced q} is a permutation between two lists of sixteen-tuples, both read from \bndc{q}
\emph{except} the coordinates the instance holds as \bndc{Coord.const} and \bndc{Coord.known}
(\bndc{coordCell}, \cref{sec:gen-boundary-m3holds}). Whether the built rows' messages are the
instance's boundary tuples is therefore decided by what \bndc{leanIsaInstance prog s} puts in
\bndc{boundary}:
\begin{itemize}
\item If the bytecode blocks are boundary blocks with
  \bndc{Coord.known (bytecodeSlotColumn prog k)} and \bndc{Coord.known (idxColumn prog.logSize)}
  (decision 8; the shape leanVM has, \bndc{layout.rs:383-395}, \bndc{Coord::Public},
  \bndc{Coord::Index}, and which the verifier evaluates itself, \bndc{leaf.rs:444, 453}), then
  for every slot \bndc{i} the instance's pushed tuple is \bndc{(g², g^i, 1, e[0], …, e[7], 0, …)}
  with \bndc{e = encodeSlots (prog.code i)} at slots 3--10 \bndc{= entry (prog.code i)}
  (\bndc{Bytecode.lean:94-96}), exactly the bus tuple of the built row's message. The same for the
  memory blocks (\bndc{known idxColumn}, \bndc{committed} limbs and counts read from \bndc{q},
  which \bndc{witnessOf} also reads) and for the state boundary (\bndc{const prog.finalPc} against
  \bndc{Boundary.lean:314}). Then \bndc{I.Balanced q}, filtered by separator, \emph{is} the three
  balances of the mixed witness, and the proof of \bndc{satisfiedBy_witnessOf} goes through.
\item If instead the two blocks were tables of the instance with flushes, obtained by
  \bndc{Component.toM3 memTable} and \bndc{Component.toM3 bytecodeTable} as the blueprint's phrase
  \enquote{\bndc{Ensemble.toM3} of the eight tables with leanISA's separators and directions}
  (\bndc{bp:803-805}) says, then their \bndc{idx}, \bndc{opcode} and \bndc{op} columns would be
  columns of \bndc{q} (the row's cells, \bndc{BytecodeRow}, \bndc{Boundary.lean:211-220}),
  \bndc{I.Balanced q} would speak of whatever the prover committed there, and \bndc{witnessOf},
  which must build the block's rows from \bndc{prog} (\bndc{bytecode_rows} demands it), would
  produce a witness whose \bndc{bytecode_balanced} does not follow from \bndc{I.Balanced q}. The
  theorem \bndc{satisfiedBy_witnessOf} would be unprovable; worse, the \emph{protocol} would be a
  different one from leanVM's (a stack with committed index and program columns, \bndc{μ} and the
  layout changed), so \bndc{verify_iff_compiled} and the Rust fixture would fail too.
\end{itemize}

\paragraph{The model.} Probe \bndc{KnownColumn} builds the spine's toy instance in both shapes:
with the block's column \bndc{known} and equal to the \enquote{program}, the honest stack
satisfies \bndc{M3Holds} and fails \bndc{Balanced} for another program, so the relation is about
the program in the instance; with the block's column \bndc{committed}, a stack that has nothing
to do with the program satisfies \bndc{M3Holds}, while the witness an adaptor would rebuild from
it, the block built from the program, does not balance. It passed at the old pins and, rewritten
with \bndc{K.ofBits 2} for the numeral \bndc{2} (which means \bndc{0} at the new CompPoly pin), as
\bndc{KnownColumnNew} at the new pins (\cref{sec:gen-boundary-probe-knowncolumn}).

\paragraph{Where \enquote{the accepted proof concerns the intended program} is decided.}
Nowhere in the oracle protocol, whose theorems hold for every \bndc{I}; in the adaptor, by the
\emph{statement} of \bndc{satisfiedBy_witnessOf}, which names \bndc{leanIsaInstance prog s} and
\bndc{SatisfiedBy prog} with one \bndc{prog}, so that its proof can only close if the instance's
\bndc{known} coordinates are the program's and the verifier's sentinel is \bndc{prog.finalPc};
and in Layer 12's \bndc{verify prog input proof}, which must evaluate those same \bndc{known}
columns (\bndc{settleFixedClaims}) and seed the transcript with the program (\bndc{fs_seed},
\bndc{cpu/mod.rs:82-93}), tied to the oracle verifier of \bndc{leanIsaInstance prog s} by
\bndc{verify_iff_compiled}. The three are consistent when they name the same \bndc{prog}. What is
missing is the lemma that lets the adaptor's proof close: Layer 2's \bndc{toM3_flushes_eq} is
about a table's flushes and cannot say anything about a boundary block. A \textbf{third bridge
lemma} is needed: for the two Clean block components and the verifier component, the messages of
the block's rows (\bndc{Boundary.lean:157-160, 240-243, 311-314}) equal, as sixteen-tuples, the
instance's boundary tuples, given the three fixed-column facts of the witness
(\cref{f:boundary-not-ensemble-tom3,f:boundary-bridge-lemma}).

\paragraph{A remark on \texorpdfstring{\protect\code{SatisfiedBy}}{SatisfiedBy} that helps.}
Its \bndc{word0_eq}, \bndc{word1_eq}, \bndc{caps.minLogMem_le}, \bndc{caps.le_maxLogMem} and
\bndc{seed_rows} are all stated on \bndc{imageOf w.data}, the \bndc{"mem"} table of the prover
data, not on the block's rows; and \bndc{MemPull.Guarantees} (\bndc{Channels.lean:214-216}) too.
\bndc{witnessOf} must therefore write the same limbs into the data and into the memory block,
and \bndc{seed_rows} is what ties them; Clean pull request \#446 would derive the data from the
block's columns and delete the conjunct (\bndc{Statement.lean:87-100}). Nothing in the chain
reads the data from anywhere else.

\subsection{The completeness direction}\label{sec:gen-boundary-completeness}

\bndc{stackOf gen w hs : Column I.μ} stacks the committed columns of \bndc{w} at the aligned
offsets (Layer 1's \bndc{Blocks.stackColumn}, \bndc{Stack.lean:70}), with \bndc{q_flock} computed
by the Flock witness generator from the \bndc{BLAKE2S} rows' limbs. \Cref{tab:gen-boundary-completeness}
goes clause by clause, from \bndc{h : SatisfiedBy prog input w} and
\bndc{hs : Sizes.ofWitness w = some s}.

{\footnotesize
\begin{longtable}{@{}>{\raggedright\arraybackslash}p{0.17\linewidth}>{\raggedright\arraybackslash}p{0.21\linewidth}>{\raggedright\arraybackslash}p{0.56\linewidth}@{}}
\caption{The clauses of \code{M3Holds} in the completeness direction \code{m3Holds_stackOf}.}\label{tab:gen-boundary-completeness}\\
\toprule
Clause & From & What else is needed \\
\midrule
\endfirsthead
\toprule
Clause & From & What else is needed \\
\midrule
\endhead
\bottomrule
\endfoot
\bndc{ConstraintsVanish} & \bndc{h.constraints}, tables 0--5 & \bndc{toM3_constraints_iff} on rows read back from the stack: \bndc{Blocks.readColumn_stackColumn} (\bndc{Stack.lean:85-87}) returns the block, so the rows read back are the rows stacked, provided every row has exactly \bndc{component.width} cells or \bndc{stackOf} reads \bndc{row[j]?.getD 0} as \bndc{Environment.fromArray} does; the \bndc{BLAKE2S} limbs read back from the strided slots equal the rows' limbs, which is the generator's \enquote{keeping the limbs in their slots} \\
\bndc{Balanced} & \bndc{h.state_balanced}, \bndc{h.mem_balanced}, \bndc{h.bytecode_balanced} & the missing boundary lemma in the other direction: the messages of \bndc{w}'s tables 6, 7 and of the verifier row are the instance's boundary tuples, which needs \bndc{h.index_columns} (the rows' \bndc{idx} cells are \bndc{g^i}, the stack does not hold them), \bndc{h.seed_rows} (the block's limbs are the data's, which the stack holds) and \bndc{h.bytecode_rows} (the rows' entry cells are \bndc{prog}'s, not stacked); then the three permutations concatenate into one, since every flush tuple carries exactly one of the three separators \\
\bndc{CountsNonzero} & \bndc{h.counts_nonzero} & \bndc{I.counts j} $⊆$ the columns that appear as a pull's count coordinate (the converse inclusion the soundness direction needs; so \bndc{counts j} must be exactly that set) \\
\bndc{PublicLinesHold} & \bndc{h.word0_eq}, \bndc{h.word1_eq}, \bndc{h.seed_rows}, \bndc{h.caps} & \bndc{MemImage.read_gpow} at \bndc{κ < 64}, \bndc{limb_ofLimbs}; the third line's \bndc{(0, 0)} from the top limb of \bndc{input.word0}, \bndc{input.word1}, which is \bndc{0} by definition (\bndc{Memory.lean:107-110}) \\
\bndc{aux} & \bndc{h.blake2s_valid} & \bndc{FlockWitnessGen}'s lemma: the wires the generator computes from limbs that compress satisfy the R1CS (decision 12) \\
\end{longtable}}

What else \bndc{m3Holds_stackOf} needs: nothing on the caps beyond
\bndc{Sizes.ofWitness w = some s}, which \bndc{h.caps.heights} (every height a power of two at
most \bndc{2^32}) supplies for the eight tables, together with \bndc{h.seed_rows} (table 6 has
\bndc{2^(imageOf w.data).1} rows) and \bndc{h.bytecode_rows} (table 7 has
\bndc{2^prog.logSize}); and \bndc{leanIsaInstance}'s layout fit, which is a \bndc{Blocks} fact
and must precede the instance (the Layer 1 dossier's finding that \bndc{leanIsaInstance_fits}
cannot be stated as written, \file{code-layer1.md}, G.2). What \bndc{verify}'s acceptance needs
beyond it is not in \bndc{SatisfiedBy} at all: the stacking window \bndc{μ ∈ [15, 28]} and the
rate window (\cref{sec:gen-boundary-hypotheses,f:boundary-base-completeness}).
